\chapter{المتسلسلات العددية}\label{ch:b1:series}

جمع أعداد لا تُحصى يعني أخذ نهاية المجاميع
الجزئية --- لا أكثر ولا أقلّ. ويهيّئ هذا الفصل
التعاريف ومحكات التقارب الصالحة للاستعمال في السنة الأولى:
المقارنة والمكافئات من أجل الحدود الموجبة، \omterm{thm:b1:series:ratio}{ومحك النسبة}، و
المقارنة التكاملية التي تعطي متسلسلات ريمان، \omterm{thm:b1:series:absolute}{والتقارب بإطلاق}،
ومبرهنة المتسلسلات المتناوبة. وأمّا النظرية الأدقّ (جداءات
المتسلسلات، والجمع بالحزم، ومتسلسلات الدوال) فتنتمي إلى
السنة الثانية.

\section{عموميات}

\begin{definition}\label{def:b1:series:def}
من أجل متتالية $(u_n)$، تكون \emph{المتسلسلة}\index{متسلسلة} $\sum u_n$
هي متتالية \emph{المجاميع الجزئية} $S_N = \sum_{n=0}^{N} u_n$.
و\emph{تتقارب} المتسلسلة إذا تقاربت $(S_N)$؛ وتكون النهاية هي
\emph{المجموع} $\sum_{n=0}^{\infty} u_n$، ويكون $R_N = \sum_{n > N} u_n =
S - S_N$ هو \emph{الباقي}، وهو يؤول إلى $0$.
\end{definition}

\begin{example}[المتسلسلة الهندسية]\label{ex:b1:series:geometric}
من أجل $q \in \C$: $\;S_N = \sum_{n=0}^{N} q^n = \frac{1 -
q^{N+1}}{1-q}$ ($q \neq 1$). وتتقارب \omterm{def:b1:series:def}{المتسلسلة} إذا وفقط إذا كان $\abs q < 1$
(\cref{exo:b1:seq:3})، مع\index{متسلسلة هندسية}
\[
\sum_{n=0}^{\infty} q^n = \frac{1}{1 - q} .
\]
\end{example}

\begin{example}[الأعداد العشرية الدورية متسلسلات هندسية]\label{ex:b1:series:periodicdecimal}
ما العدد $0.363636\dots$؟ إن كتابته نفسها \omterm{def:b1:series:def}{متسلسلة}:
\[
0.\overline{36} = \sum_{k=1}^{\infty} \frac{36}{100^k}
= 36\cdot\frac{1/100}{1 - 1/100} = \frac{36}{99} =
\frac{4}{11} ,
\]
بالمجموع الهندسي مع $q = \frac{1}{100}$. وعمومًا فإن
كتلة $B$ من $p$ رقمًا تتكرر إلى الأبد تساوي
$\frac{B}{10^p - 1}$ --- وهي الآلية وراء محك الدورية
في \cref{pb:b1:reals:1}، الذي تذكره لغة هذا الفصل
أخيرًا في سطر واحد: \omterm{pb:b1:reals:1}{فالنشر العشري} \omterm{def:b1:series:def}{متسلسلة}
متقاربة، دورية في النهاية إذا وفقط إذا كان مجموعها ناطقًا.
وقد كانت آلة الأرقام في الفصل 10، المبنية هناك بالحدود
العليا العارية، نظريةَ متسلسلات مسافرةً متنكّرة.
\end{example}

\begin{proposition}[الوقائع الأولى]\label{prop:b1:series:first}
\begin{enumerate}
  \item إذا تقاربت $\sum u_n$، فإن $u_n \to 0$. (والعكس
        \emph{خاطئ}: \omterm{def:b1:series:def}{فالمتسلسلة} التوافقية.)
  \item الخطية: تُجمع المتسلسلات المتقاربة وتُضرب في سلّم، بالمجاميع
        المتوقّعة.
  \item (التلسكوب\index{متسلسلة تلسكوبية}) تتقارب $\sum (v_{n+1} -
        v_n)$ إذا وفقط إذا تقاربت $(v_n)$، بالمجموع $\lim v_n -
        v_0$.
  \item تغيير عدد منته من الحدود لا يؤثر في التقارب
        (بل في المجموع فقط).
\end{enumerate}
\end{proposition}

\begin{proof}
(1) $u_N = S_N - S_{N-1} \to S - S = 0$. وللمتسلسلة التوافقية
$u_n = \frac1n \to 0$ ومع ذلك تتباعد (\cref{exo:b1:seq:5}). (2)
عمليات على النهايات. (3) $S_N = v_{N+1} - v_0$. (4) تتغير المجاميع
الجزئية بمقدار ثابت في النهاية.
\end{proof}

\begin{example}[تخطيط الأرقام بالباقي الهندسي]\label{ex:b1:series:georemainder}
من أجل $\abs q < 1$ يكون باقي \omterm{ex:b1:series:geometric}{المتسلسلة الهندسية}
صريحًا:
\[
R_N = \sum_{n = N+1}^{\infty} q^n = \frac{q^{N+1}}{1 - q} .
\]
وهذا يحوّل أهداف الدقة إلى أعداد حدود قبل أيّ
حساب. فلتقويم $\sum_{n\geq0} \bigl(\frac13\bigr)^n =
\frac32$ بدقة $10^{-10}$: نحتاج إلى $\frac{(1/3)^{N+1}}{2/3} \leq
10^{-10}$، أي $3^{N} \geq \frac{3}{2}\cdot 10^{10}$، أي $N
\geq 22$ (لأن $3^{22} \approx 3.1\cdot10^{10}$): أي ثلاثة وعشرون
حدًّا، معلومة مسبقًا. وكل تقدير بسرعة هندسية في
مسائل نهاية الأسبوع (\omterm{def:b1:series:def}{متسلسلة} $\frac13$ من أجل $\ln 2$، وظلال ماشان
العكسية في \cref{pb:b1:taylor:1}) هو هذه الميزانية ذات السطرين
في ثوب احترافي.
\end{example}

\begin{example}[تلسكوب أطول]\label{ex:b1:series:telescope3}
احسب $\sum_{n\geq1} \frac{1}{n(n+1)(n+2)}$. بالعناصر
البسيطة (\cref{ch:b1:fractions}):
\[
\frac{1}{n(n+1)(n+2)}
= \frac{1/2}{n} - \frac{1}{n+1} + \frac{1/2}{n+2}
= \frac12\Bigl(\frac{1}{n(n+1)} - \frac{1}{(n+1)(n+2)}\Bigr),
\]
حيث تكون الصورة الثانية --- وهي فرق قيمتين متتاليتين للمقدار
$w_n = \frac{1}{n(n+1)}$ --- هي التلسكوبية. ومنه
\[
\sum_{n=1}^{N} \frac{1}{n(n+1)(n+2)}
= \frac12\Bigl(w_1 - w_{N+1}\Bigr)
= \frac12\Bigl(\frac12 - \frac{1}{(N+1)(N+2)}\Bigr)
\longrightarrow \frac14 .
\]
والفكرة النافذة: قلّما تتلسكب العناصر البسيطة الثلاثية
كما تُكتب؛ فأعد تجميعها أولًا في فرق $w_n -
w_{n+1}$ --- والجائزة ليست التقارب فحسب بل
المجموع المضبوط، وهو ما لا يعطيه أيّ محك مقارنة.
\end{example}

\section{المتسلسلات ذات الحدود غير السالبة}

\begin{theorem}[المجاميع الجزئية المحدودة]\label{thm:b1:series:positive}
إذا كان $u_n \geq 0$ من أجل كل $n$، تتزايد المجاميع الجزئية، ومنه: تتقارب $\sum
u_n$ $\iff$ كانت مجاميعها الجزئية محدودة من أعلى. ومنه
\emph{محك المقارنة}: إذا كان $0 \leq u_n \leq v_n$ من أجل كل $n$ (كبير)،
\[
\sum v_n \text{ تتقارب} \implies \sum u_n \text{ تتقارب},
\qquad
\sum u_n \text{ تتباعد} \implies \sum v_n \text{ تتباعد}.
\]
و\emph{محك المكافئات}: إذا كان $u_n \sim v_n$ مع $v_n \geq
0$، فللمتسلسلتين الطبيعة نفسها.
\end{theorem}

\begin{proof}
مبرهنة النهاية الرتيبة (\cref{thm:b1:seq:monotone}) من أجل النقطة
الأولى؛ ومقارنة المجاميع الجزئية من أجل الثانية. وأمّا المكافئات: فمن أجل
$n$ كبير، $\frac12 v_n \leq u_n \leq 2 v_n$ (وهو تعريف $\sim$
مع $\varepsilon = \frac12$)، وتنطبق المقارنة في الاتجاهين.
\end{proof}

\begin{example}[مكافئ يبرهن على التباعد]\label{ex:b1:series:equivdiverge}
ما طبيعة $\sum_{n\geq1} n\sin\dfrac{1}{n^2}$؟ لأن
$\frac{1}{n^2} \to 0$ و $\sin h \sim h$ عند $0$:
\[
n\sin\frac{1}{n^2} \;\sim\; n\cdot\frac{1}{n^2} = \frac1n ,
\]
وينقل محك المكافئات تباعدَ
\omterm{def:b1:series:def}{المتسلسلة} التوافقية: فهي متباعدة --- وإن آلت حدودها إلى
$0$. أي نشر واحد وسلّم واحد وحكم واحد؛ والنمط
نفسه ذو الخطوتين (مكافئ، ثم بحث في ريمان أو الهندسية) يفصل
في متسلسلات \cref{exo:b1:series:3} الأربع كلها.
\end{example}

\begin{example}[محك المكافئات في سطر واحد]\label{ex:b1:series:equivtest}
ما طبيعة $\sum_{n \geq 1} \frac{\sqrt{n+1} - \sqrt n}{n}$؟
رافِق البسط:
\[
\frac{\sqrt{n+1} - \sqrt n}{n}
= \frac{1}{n\,(\sqrt{n+1} + \sqrt n)}
\sim \frac{1}{2\,n^{3/2}} ,
\]
وهو سلّم ريمان متقارب ($\alpha = \frac32 > 1$): فتتقارب \omterm{def:b1:series:def}{المتسلسلة}.
وقد أخذ القرار كله مكافئًا واحدًا وبحثًا واحدًا
--- شرط أن تكون الحدود غير سالبة، وهي كذلك. والفكرة
النافذة: من أجل المتسلسلات الموجبة، تكون نظرية التقارب كلها
\emph{قاموسَ سلالم} ($n^{-\alpha}$ و $q^n$ و
$\frac{1}{n(\ln n)^\alpha}$) مع رخصة استبدال حدّ
بمكافئه؛ والعمل التحليلي في المقاربات
(\cref{ch:b1:taylor})، لا في الجمع أبدًا.
\end{example}

\begin{theorem}[المقارنة التكاملية؛ متسلسلات ريمان]\label{thm:b1:series:riemann}
لتكن $f$ \omterm{def:b1:continuity:continuous}{متصلة} وغير سالبة و\emph{متناقصة} على
$\intco{1}{+\infty}$. عندئذ
\[
\int_1^{N+1} f(t)\,\dd t \;\leq\; \sum_{n=1}^{N} f(n) \;\leq\; f(1)
+ \int_1^{N} f(t)\,\dd t ,
\]
ومنه تتقارب $\sum f(n)$ إذا وفقط إذا كان $\bigl(\int_1^x f\bigr)$ محدودًا. وعلى
الخصوص، من أجل $\alpha \in \R$:\index{متسلسلة ريمان}
\[
\sum_{n \geq 1} \frac{1}{n^\alpha} \text{ تتقارب}
\iff \alpha > 1,
\]
و $\sum_{n=1}^{N} \frac1n = \ln N + O(1)$.
\end{theorem}

\begin{proof}
من أجل $n \leq t \leq n+1$، تعطي الرتابة أن $f(n+1) \leq f(t)
\leq f(n)$؛ وبالمكاملة على $\intcc{n}{n+1}$ (وهي قطعة
طولها $1$):
\[
f(n+1) \;\leq\; \int_n^{n+1} f(t)\,\dd t \;\leq\; f(n) .
\]
وبجمع المتراجحات اليمنى من أجل $n = 1, \dots, N-1$ نجد
$\int_1^{N} f \leq \sum_{n=1}^{N-1} f(n)$، ومنه التأطير
الأعلى بعد إضافة $f(N) \leq f(1)$؛ وبجمع اليسرى من أجل
$n = 1, \dots, N$ نجد $\sum_{n=2}^{N+1} f(n) \leq
\int_1^{N+1} f$، وهو بعد إعادة التدليل التأطير الأدنى.
والتقارب: تحدّ المجاميع الجزئية والتكاملات $\int_1^x f$
بعضها بعضًا في حدود الثابت $f(1)$، وكلاهما
غير متناقص، ومنه يكون أحدهما محدودًا إذا وفقط إذا كان الآخر كذلك
(\cref{thm:b1:series:positive}). ومن أجل $f(t)
= t^{-\alpha}$ ($\alpha \neq 1$): $\int_1^x t^{-\alpha}\dd t =
\frac{x^{1-\alpha} - 1}{1 - \alpha}$، وهو محدود إذا وفقط إذا كان $\alpha > 1$؛ ومن أجل
$\alpha = 1$ يكون \omterm{thm:b1:integration:def}{التكامل} $\ln x \to \infty$، ويعطي التأطير
$\ln(N+1) \leq H_N \leq 1 + \ln N$. ومن أجل $\alpha \leq 0$ لا تؤول
الحدود إلى $0$.
\end{proof}

\begin{example}[الكومة التوافقية]\label{ex:b1:series:harmonicstack}
كم حدًّا يجب أن تراكم \omterm{def:b1:series:def}{المتسلسلة} التوافقية لتتجاوز $20$؟
يجيب التأطير $\ln(N+1) \leq H_N \leq 1 + \ln N$ دون أيّ
جمع أصلًا: فالشرط $H_N \geq 20$ يقتضي $1 + \ln N \geq 20$،
أي $N \geq \eu^{19} \approx 1.8\cdot10^{8}$، وهو مضمون
ما إن يكون $\ln(N + 1) \geq 20$، أي $N \approx \eu^{20} \approx
4.9\cdot10^{8}$. (وتلطّف مسألة نهاية الأسبوع هذا إلى $N \approx
\eu^{20 - \gamma} \approx 2.7\cdot10^{8}$ عبر \omterm{pb:b1:series:1}{ثابت أويلر}.)
والفكرة النافذة: لا تكتفي المقارنة التكاملية بالفصل
في التقارب --- بل \emph{تحدّد موضع} المجاميع الجزئية بدقة
لوغاريتمية، فتحوّل حسابًا ميؤوسًا منه
(مئات الملايين من الحدود) إلى تقدير في سطرين.
\end{example}

\begin{theorem}[محك النسبة (دالامبير)]\label{thm:b1:series:ratio}
لتكن $u_n > 0$ مع $\frac{u_{n+1}}{u_n} \to \ell$.\index{محك النسبة}
\begin{itemize}
  \item إذا كان $\ell < 1$: تتقارب $\sum u_n$؛
  \item وإذا كان $\ell > 1$: $u_n \to +\infty$، أي التباعد؛
  \item وإذا كان $\ell = 1$: فلا نتيجة (\omterm{def:b1:series:def}{فالمتسلسلة} $\sum \frac1n$ تتباعد و $\sum
        \frac{1}{n^2}$ تتقارب).
\end{itemize}
\end{theorem}

\begin{proof}
إذا كان $\ell < 1$، فثبّت $q \in \intoo{\ell}{1}$: فبعد $N$ ما،
$u_{n+1} \leq q\,u_n$، ومنه $u_n \leq u_N q^{\,n-N}$ بالاستقراء:
أي المقارنة \omterm{ex:b1:series:geometric}{بمتسلسلة هندسية}. وإذا كان $\ell > 1$: فبعد $N$ ما
تكون المتتالية $(u_n)$ متزايدة، ومنه لا يمكن أن تؤول إلى $0$
(لأن نهايتها، إن وُجدت، هي $\geq u_N > 0$)؛ ومنه، حسب
\cref{prop:b1:series:first}~(1)، يقع التباعد --- بل إن $u_n
\geq u_N q^{n-N}$ مع $q > 1$ يعطي $u_n \to \infty$.
\end{proof}

\begin{example}\label{ex:b1:series:ratio}
تتقارب $\sum \frac{x^n}{n!}$ من أجل كل $x > 0$: فالنسبة
$\frac{x}{n+1} \to 0$. ومجموعها هو $\eu^x$: فبتايلور--لاغرانج
(\cref{thm:b1:taylor:lagrange}) على $\intcc{0}{x}$،
\[
\Bigl| \eu^x - \sum_{k=0}^{n} \frac{x^k}{k!} \Bigr|
\leq \eu^{x}\, \frac{x^{n+1}}{(n+1)!} \xrightarrow[n\to\infty]{} 0 ,
\]
ويؤول الحاصر إلى $0$ لأن العاملي يهيمن
(\cref{exo:b1:integration:9}~(1) استعمل الواقعة نفسها). والحجة نفسها
تجمع متسلسلات $\sin$ و $\cos$ و $\sinh$ و $\cosh$ على $\R$ كله.
\end{example}

\begin{example}[محك النسبة كافٍ لا لازم]\label{ex:b1:series:rationotnecessary}
لتكن $u_n = 2^{-n}$ من أجل $n$ زوجي و $u_n = 2^{-n-2}$ من أجل $n$
فردي. فتتذبذب النسب المتتالية بين $\frac{1}{8}$ و
$\frac12\cdot4 = 2$، ومنه فلا نهاية للمقدار $\frac{u_{n+1}}{u_n}$ و
يصمت دالامبير --- ومع ذلك $u_n \leq 2^{-n}$ ويفصل محك
المقارنة في التقارب فورًا. ففرض المحك (أي أن
\emph{تتقارب} النسبة) قيدٌ حقيقي: فهو يناسب الحدود
ذات البنية الضربية المهيمنة الواحدة (العامليات والقوى)،
ويفشل على كل ما يتنفس. وحين تسيء النسب السلوك،
فتراجع إلى المقارنة بغلاف هندسي --- وهو كل ما كان
\omterm{thm:b1:series:ratio}{محك النسبة} أصلًا، كما يبيّن برهانه.
\end{example}

\begin{example}[محك النسبة في معارك العامليات]\label{ex:b1:series:ratiobinom}
ما طبيعة $\sum_{n\geq0} \dfrac{(n!)^2}{(2n)!}$ (مقلوبات
المعاملات الثنائية المركزية، مضروبةً في العامل $n + 1$)؟
تُسقط النسبة العامليات:
\[
\frac{u_{n+1}}{u_n}
= \frac{((n+1)!)^2}{(n!)^2}\cdot\frac{(2n)!}{(2n+2)!}
= \frac{(n+1)^2}{(2n+1)(2n+2)}
\longrightarrow \frac14 < 1 :
\]
فهي متقاربة، مع فسحة --- إذ تتلاشى الحدود جوهريًا
مثل $4^{-n}$، وهذا متسق مع $\binom{2n}{n} \geq
\frac{4^n}{2n+1}$ من \cref{pb:b1:integration:1}. والفكرة
النافذة: قسمة العامليات هي بالضبط ما يهضمه \omterm{thm:b1:series:ratio}{محك النسبة}
--- فكل عاملي يختصر إلى \omterm{def:b1:fractions:field}{كسر ناطق} في
$n$، تُقرأ نهايته من الحدود المهيمنة.
\end{example}

\section{التقارب بإطلاق؛ المتسلسلات المتناوبة}

\begin{theorem}[التقارب بإطلاق]\label{thm:b1:series:absolute}
إذا تقاربت $\sum \abs{u_n}$ (أي \emph{التقارب
بإطلاق}\index{تقارب بإطلاق})، فإن $\sum u_n$ تتقارب،
و $\bigl|\sum u_n\bigr| \leq \sum \abs{u_n}$. وهذا يصحّ من أجل الحدود الحقيقية
أو العقدية.
\end{theorem}

\begin{proof}
تحقق المجاميع الجزئية، من أجل $M > N$ (بمحك كوشي،
\cref{thm:b1:seq:complete}):
\[
\abs{S_M - S_N} = \Bigl| \sum_{n=N+1}^{M} u_n \Bigr|
\leq \sum_{n=N+1}^{M} \abs{u_n},
\]
وهو صغير من أجل $N$ كبير لأن المجاميع الجزئية للمقدار $\sum\abs{u_n}$
تكوّن \omterm{def:b1:seq:cauchy}{متتالية كوشي}. ومنه فالمتتالية $(S_N)$ \omterm{def:b1:seq:cauchy}{متتالية كوشي}، ومنه متقاربة. و
تمرّ المتراجحة إلى النهاية من متراجحة المثلث المنتهية.
\end{proof}

\begin{example}[التقارب بإطلاق، حقيقيًا وعقديًا]\label{ex:b1:series:absexamples}
$\sum_{n\geq1} \frac{\sin n}{n^2}$: تتغير إشارات الحدود
تغيّرًا مضطربًا (بل إن $(\sin n)$ \omterm{def:b1:topology:dense}{كثيفة} في $\intcc{-1}{1}$،
\cref{exo:b1:seq:12})، ولا بنية تناوب تلوح
في الأفق. وينقذ \omterm{thm:b1:series:absolute}{التقارب بإطلاق} كل شيء دفعة واحدة:
$\bigl|\frac{\sin n}{n^2}\bigr| \leq \frac{1}{n^2}$، وهو
سلّم متقارب، ومنه تتقارب \omterm{def:b1:series:def}{المتسلسلة}. ويعمل الدرع نفسه
على $\C$: فتتقارب $\sum_{n\geq1}\frac{\eu^{\iu n}}{n^2}$
لأن $\bigl|\frac{\eu^{\iu n}}{n^2}\bigr| = \frac{1}{n^2}$
--- فأنماط الإشارة، حتى الثنائية البعد منها، لا شأن لها
ما إن تكون المقاييس قابلة للجمع. والفكرة النافذة: \omterm{thm:b1:series:absolute}{التقارب
بإطلاق} هو الأداة الوحيدة في هذا الفصل التي لا تسأل أبدًا كيف
تنظَّم الإشارات؛ فجرّبه أولًا
(\cref{met:b1:series:decide})، واحفظ المحكات الرقيقة
للمتسلسلات التي تفشل فيه.
\end{example}

\begin{theorem}[محك المتسلسلات المتناوبة]\label{thm:b1:series:alternating}
لتكن $(a_n)$ متناقصة مع $a_n \to 0$. عندئذ تتقارب \omterm{thm:b1:series:alternating}{المتسلسلة
المتناوبة} $\sum (-1)^n a_n$؛ ويقع مجموعها بين أيّ مجموعين جزئيين
متتاليين، و\index{متسلسلة متناوبة}
\[
\abs{R_N} = \Bigl| \sum_{n > N} (-1)^n a_n \Bigr| \leq a_{N+1} .
\]
\end{theorem}

\begin{proof}
المجاميع الجزئية الزوجية والفردية متجاورة: $S_{2p+2} - S_{2p} =
a_{2p+2} - a_{2p+1} \leq 0$ (متناقصة) و $S_{2p+1} - S_{2p-1} =
a_{2p} - a_{2p+1} \geq 0$ (متزايدة) و $S_{2p} - S_{2p+1} =
a_{2p+1} \to 0$. وحسب \cref{thm:b1:seq:adjacent} تتقاسمان نهاية
$S$، ويجعلها محك المتتاليتين الجزئيتين
(\cref{prop:b1:seq:subsequences}) نهايةً للمتتالية $(S_N)$؛
وفوق ذلك يُحبس $S$ بين مجموعين جزئيين متتاليين، ويكون
$\abs{S - S_N}$ على الأكثر الفجوةَ إلى التالي، أي $a_{N+1}$.
\end{proof}

\begin{example}[المتسلسلة التوافقية المتناوبة]\label{ex:b1:series:ln2}
تتقارب $\sum_{n \geq 1} \frac{(-1)^{n-1}}{n}$ (بالمحك المتناوب)
لكن لا بإطلاق (لأن \omterm{def:b1:series:def}{المتسلسلة} التوافقية). ومجموعها هو $\ln 2$: فمن
المتطابقة الهندسية المنتهية $\frac{1}{1+t} = \sum_{k=0}^{n-1} (-t)^k +
\frac{(-t)^n}{1+t}$، بالمكاملة على $\intcc{0}{1}$:
\[
\ln 2 = \sum_{k=1}^{n} \frac{(-1)^{k-1}}{k}
+ (-1)^n \int_0^1 \frac{t^n}{1+t}\,\dd t ,
\qquad
0 \leq \int_0^1 \frac{t^n}{1+t}\,\dd t \leq \frac{1}{n+1} \to 0 .
\]
والتقارب بطيء بشكل مؤلم ($R_N \approx \frac{1}{N}$) ---
فالمتسلسلات المتناوبة تتقارب بالتلاشي، لا بالصغر.
\end{example}

\begin{omfigure}
\begin{tikzpicture}[x=0.78cm, y=7.2cm]
  \draw[-Stealth] (0.3,0.42) -- (11,0.42)
    node[below, font=\small] {$N$};
  \draw (1,0.415) -- (1,0.425) node[below=1pt, font=\small]
    {$1$};
  \draw (5,0.415) -- (5,0.425) node[below=1pt, font=\small]
    {$5$};
  \draw (10,0.415) -- (10,0.425) node[below=1pt, font=\small]
    {$10$};
  \draw[gray, dashed] (0.3,0.6931) -- (10.7,0.6931)
    node[right, font=\small, text=black] {$\ln 2$};
  \node[left, font=\small] at (0.3,1) {$1$};
  \node[left, font=\small] at (0.3,0.5) {$0.5$};
  \draw[omDef, thick]
    plot coordinates {(1,1) (2,0.5) (3,0.8333) (4,0.5833)
    (5,0.7833) (6,0.6167) (7,0.7595) (8,0.6345)
    (9,0.7456) (10,0.6456)};
  \fill[omDef] (1,1) circle (1.5pt);
  \fill[omDef] (2,0.5) circle (1.5pt);
  \fill[omDef] (3,0.8333) circle (1.5pt);
  \fill[omDef] (4,0.5833) circle (1.5pt);
  \fill[omDef] (5,0.7833) circle (1.5pt);
  \fill[omDef] (6,0.6167) circle (1.5pt);
  \fill[omDef] (7,0.7595) circle (1.5pt);
  \fill[omDef] (8,0.6345) circle (1.5pt);
  \fill[omDef] (9,0.7456) circle (1.5pt);
  \fill[omDef] (10,0.6456) circle (1.5pt);
\end{tikzpicture}

{\small المجاميع الجزئية $S_N$ \omterm{def:b1:series:def}{للمتسلسلة} التوافقية
المتناوبة $1 - \frac12 + \frac13 - \cdots$ تقفز فوق نهايتها
$\ln 2$ عند كل خطوة: فالمجاميع الفردية من أعلى، والزوجية من
أسفل، وكل قفزة حجمها $\frac{1}{N+1}$. والتأطير هو
برهان \cref{thm:b1:series:alternating} مجعولًا مرئيًا --- و
الانغلاق البطيء للكمّاشة ($\abs{S_N - \ln 2} \approx
\frac{1}{2N}$، ومسألة نهاية الأسبوع \cref{pb:b1:series:1}) هو سبب أن
لا أحد يحسب $\ln 2$ بهذه الكيفية.}
\end{omfigure}

\begin{remark}[مزالق شائعة مع المتسلسلات]
(أ) \emph{يحتاج محك المكافئات إلى إشارة}: لتكن $v_n =
\frac{(-1)^n}{\sqrt n}$ و $u_n = v_n + \frac1n$. عندئذ
$\frac{u_n}{v_n} = 1 + \frac{(-1)^n}{\sqrt n} \to 1$، ومنه $u_n
\sim v_n$؛ ومع ذلك تتقارب $\sum v_n$ (بالمحك المتناوب) بينما
تتباعد $\sum u_n = \sum v_n + \sum \frac1n$. فالتكافؤ
يضبط \emph{حجم} الحدود، وأمّا في المتسلسلات ذات الإشارة فالحجم
ليس قدرًا --- والمحك مذكور، وصحيح، من أجل الحدود
غير السالبة (في النهاية) فقط. (ب) \emph{المقدار $u_n \to 0$
لا يبرهن على شيء}: \omterm{def:b1:series:def}{فالمتسلسلة} التوافقية المثال المضاد
الأبدي؛ والاتجاه العكسي
(\cref{prop:b1:series:first}~(1)) ليس إلا محك تباعد
سريعًا. (ج) \emph{نهاية النسبة $1$ صمتٌ لا
تقارب}: فلكلٍّ من $\sum\frac1n$ و $\sum\frac{1}{n^2}$
نسبةٌ $\to 1$؛ فانتقل إلى سلالم ريمان أو المقارنة التكاملية.
(د) \emph{يحتاج المتناوب إلى التناقص}: فالمقدار $\sum
\frac{(-1)^n}{n + (-1)^n}$ يبدو متناوبًا ولا يُعالج
إلا بالنشر (\cref{exo:b1:series:5})؛ وتبيّن مسألة نهاية الأسبوع
في \cref{ch:b1:taylor} (السؤال 23 هناك) أن المحك قد
يفشل تمامًا دون رتابة. (ه) \emph{التجميع و
إعادة الترتيب ليسا مجانيين}: فإدخال الأقواس غير مؤذٍ في
المتسلسلات المتقاربة لكنه قد ينشئ تقاربًا من تباعد
($1 - 1 + 1 - \cdots$ مجمَّعةً أزواجًا)، وقد تغيّر إعادة الترتيب
المجموعَ نفسه --- وهي الدراما المعروضة في مسألة نهاية الأسبوع
في هذا الفصل (\cref{pb:b1:series:1}).
\end{remark}

\begin{method}[الفصل في طبيعة متسلسلة]\label{met:b1:series:decide}
\begin{enumerate}
  \item هل $u_n \to 0$؟ إن لم يكن، فالتباعد، وتوقّف.
  \item الحدود غير السالبة: ابحث عن مكافئ للمقدار $u_n$ (بالنشور،
        \cref{ch:b1:taylor}!)، وقارن بسلالم ريمان أو
        الهندسية؛ وتستدعي العامليات والقوى \omterm{thm:b1:series:ratio}{محك النسبة}؛
        وتستدعي $f(n)$ المتناقصة المقارنةَ التكاملية.
  \item الإشارات متغيرة: جرّب \omterm{thm:b1:series:absolute}{التقارب بإطلاق} أولًا؛ فإن فشل، فالمحك
        المتناوب (وتحقق من \emph{التناقص} بعناية)؛ وما بعد
        ذلك، أدوات السنة الثانية.
\end{enumerate}
\end{method}

\begin{example}[مقامات فردية ونصف تلسكوب]\label{ex:b1:series:oddtelescope}
احسب $\sum_{n \geq 1} \dfrac{1}{4n^2 - 1}$. بالعناصر
البسيطة: $\frac{1}{(2n-1)(2n+1)} =
\frac12\bigl(\frac{1}{2n-1} - \frac{1}{2n+1}\bigr)$، ومنه
\[
\sum_{n=1}^{N} \frac{1}{4n^2 - 1}
= \frac12\Bigl(1 - \frac{1}{2N+1}\Bigr)
\longrightarrow \frac12 .
\]
وقارن بالمقدار $\sum \frac{1}{n(n+1)} = 1$
(\cref{exo:b1:series:1}): فالهيكل التلسكوبي نفسه، لكن
الحدين المتتاليين هنا يبعد أحدهما عن الآخر باثنين في الأعداد الفردية، و
يسجّل العامل $\frac12$ الخطوة. والفكرة النافذة:
التلسكوب تغييرُ منظور لا حيلة --- فكلما كان
الحدّ العام فرقًا $w_n - w_{n+1}$ لمتتالية
لها نهاية، كان المجموع $w_1 - \lim w$، وهو بالضبط
\cref{prop:b1:series:first}~(3).
\end{example}

\begin{remark}[سلسلة التحليل، بأثر رجعي]
هذا الفصل هو حيث يتقارب تحليل المجلّد، وكل
محك يسمّي سلفه. فالمجاميع الجزئية المحدودة هي مبرهنة النهاية
الرتيبة (\cref{ch:b1:seq})، وهي نفسها بديهية التمام
في \cref{ch:b1:reals}؛ \omterm{thm:b1:series:absolute}{والتقارب بإطلاق} هو محك
كوشي؛ والمحك التكاملي هو تأطير \cref{ch:b1:integration}
للمساحات؛ ومكافئات الحدود العامة هي
نشور \cref{ch:b1:taylor}؛ والمبرهنة المتناوبة هي
مبرهنة المتتاليتين المتجاورتين المساعدة في ثياب العيد. وبالقراءة
بالمقلوب، تفسّر السلسلة \emph{لماذا} وُجد كل فصل
--- ومسائل نهاية الأسبوع المنظومة عبرها
(أرقام الأساس $b$، وتشيزارو--ستولتز، وآلات الصمم،
\omterm{pb:b1:series:1}{وثابت أويلر}) هي الأفكار القليلة نفسها ملتقيةً على ارتفاع
أعلى فأعلى. والجبر الخطي الذي يليه يغيّر
الموضوع لا المعايير: فعادة العبارات المضبوطة بخطأ
مصادَق عليه تنجو من الانتقال من النهايات إلى الأبعاد.
\end{remark}

\begin{remark}[إلى أين تذهب المتسلسلات تاليًا]
يُغلق هذا الفصل تحليل المجلّد ويفتح ثلاثة
أبواب. ففي مجلّد السنة الثانية، تكتسب المتسلسلات متغيرًا ($\sum
a_n x^n$: المتسلسلات الصحيحة، بنصف قطر تقاربها) ثم
نظريةً ذات قيم دالية (متسلسلات فورييه)؛ وتصير ثنائية
\omterm{thm:b1:series:absolute}{التقارب بإطلاق} إزاء التقارب الشرطي، المعروضة دراميًا في
مسألة نهاية الأسبوع أدناه، حجرَ زاوية لكليهما. وفي
الاحتمالات (مجلّد السنة الثالثة)، تكون آمال المتغيرات العشوائية
المتقطّعة \emph{هي} متسلسلات، ويكون \omterm{thm:b1:series:absolute}{التقارب بإطلاق} هو ما
يجعلها معرَّفة جيدًا. \omterm{thm:b1:series:riemann}{ومتسلسلة ريمان} $\sum n^{-s}$،
مدفوعةً إلى $s$ عقدي، تصير دالة زيتا --- وهي
أكثر \omterm{def:b1:series:def}{متسلسلة} دُرست في الرياضيات.
\end{remark}

\section{تمارين}

\begin{exercise}[$\star$]\label{exo:b1:series:1}
الطبيعة (والمجموع، حين يقع التلسكوب) من أجل:
\[
\sum_{n\geq1} \frac{1}{n(n+1)},
\qquad
\sum_{n\geq2} \ln\Bigl(1 - \frac{1}{n^2}\Bigr),
\qquad
\sum_{n\geq0} \frac{3^n + 4^n}{5^n} .
\]
\end{exercise}

\begin{exercise}[$\star$]\label{exo:b1:series:2}
الطبيعة من أجل: $\;\sum \dfrac{n^2}{2^n}$؛ $\;\sum \dfrac{n!}{n^n}$؛
$\;\sum \dfrac{2^n\,n!}{n^n}$؛ $\;\sum \dfrac{3^n\,n!}{n^n}$.
\emph{(\omterm{thm:b1:series:ratio}{بمحك النسبة}؛ وتذكّر $\bigl(1 + \frac1n\bigr)^n \to \eu$.)}
\end{exercise}

\begin{exercise}[$\star$]\label{exo:b1:series:3}
الطبيعة من أجل: $\;\sum \sin\dfrac{1}{n^2}$؛ $\;\sum
\Bigl(1 - \cos\dfrac1n\Bigr)$؛ $\;\sum \dfrac{1}{\sqrt{n(n+1)}}$؛
$\;\sum \dfrac{\ln n}{n^2}$ \emph{(قارن بالمقدار $n^{-3/2}$)}.
\end{exercise}

\begin{exercise}[$\star$]\label{exo:b1:series:4}
برهن على أن $\sum_{n\geq1} \frac{1}{n^2}$ تتقارب بمجموع $\leq 2$،
باستعمال $\frac{1}{n^2} \leq \frac{1}{n(n-1)}$ من أجل $n \geq 2$ وحاصر
تلسكوبي.
\end{exercise}

\begin{exercise}[$\star\star$]\label{exo:b1:series:5}
الطبيعة من أجل $\;\sum \dfrac{(-1)^n}{\sqrt n}$، ومن أجل $\;\sum
\dfrac{(-1)^n}{n + (-1)^n}$ \emph{(انشر: فالمحك المتناوب لا
ينطبق مباشرة --- ولماذا؟)}، ومن أجل $\;\sum
\sin\bigl(\pi\sqrt{n^2+1}\,\bigr)$ \emph{(ردّ بترديد $\pi$:
$\sqrt{n^2+1} = n + \frac{1}{2n} + O(n^{-3})$)}.
\end{exercise}

\begin{exercise}[$\star\star$]\label{exo:b1:series:6}
من أجل أيّ $\alpha > 0$ تتقارب $\sum_{n \geq 2}
\dfrac{1}{n (\ln n)^{\alpha}}$؟ \emph{(بالمقارنة التكاملية؛ و
عوّض $u = \ln t$.)}
\end{exercise}

\begin{exercise}[$\star\star$]\label{exo:b1:series:7}
لتكن $u_n = \dfrac{1}{n} - \ln\Bigl(1 + \dfrac1n\Bigr)$. برهن على أن
$0 \leq u_n \leq \dfrac{1}{2n^2}$، وعلى أن $\sum u_n$ تتقارب، و
استنتج وجود \omterm{pb:b1:series:1}{ثابت أويلر}\index{ثابت أويلر}:
\[
\gamma = \lim_{N \to \infty} \Bigl( \sum_{n=1}^{N} \frac 1n - \ln N
\Bigr) .
\]
\end{exercise}

\begin{exercise}[$\star\star$]\label{exo:b1:series:8}
احسب المجموعين
\[
\sum_{n=1}^{\infty} \frac{1}{n(n+2)}
\qquad\text{و}\qquad
\sum_{n=0}^{\infty} \frac{n}{2^n} .
\]
\emph{(من أجل الأول: بالعناصر البسيطة. ومن أجل الثاني: احسب
$\sum_{n=1}^{N} n x^{n-1}$ في صورة \omterm{def:b1:topology:closed}{مغلقة} ودع $N \to \infty$ عند
$x = \frac12$.)}
\end{exercise}

\begin{exercise}[$\star\star\star$]\label{exo:b1:series:9}
(تكثيف كوشي) لتكن $(u_n)$ غير سالبة ومتناقصة.
برهن على أن
\[
\sum_{n \geq 1} u_n \text{ تتقارب}
\iff
\sum_{k \geq 0} 2^k\, u_{2^k} \text{ تتقارب},
\]
بمقارنة حزم الحدود بين قوى $2$ المتتالية.
واستعد منها محك ريمان و
\cref{exo:b1:series:6}.
\end{exercise}

\begin{exercise}[$\star\star\star$]\label{exo:b1:series:10}
باستعمال المتطابقة التكاملية في \cref{ex:b1:series:ln2} مكيَّفةً على
$\frac{1}{1+t^2}$، برهن على صيغة لايبنتز
\[
\frac{\pi}{4} = \sum_{n=0}^{\infty} \frac{(-1)^n}{2n+1}
= 1 - \frac13 + \frac15 - \frac17 + \cdots
\]
مع حاصر الخطأ $\abs{R_N} \leq \frac{1}{2N+3}$.
\end{exercise}

\begin{exercise}[$\star\star$]\label{exo:b1:series:11}
الطبيعة من أجل $\displaystyle\sum_{n \geq 1} \frac{1}{n^{1 + 1/n}}$.
\emph{(احسب نهاية $n^{1/n}$ وجد مكافئًا للحدّ
العام: فمحك ريمان يحتاج إلى أُسّ \emph{مثبَّت}.)}
\end{exercise}

\begin{exercise}[$\star\star\star$]\label{exo:b1:series:12}
لتكن $(u_n)$ غير سالبة و\emph{متناقصة} مع $\sum u_n$
متقاربة. برهن على أن $n\,u_n \to 0$ \emph{(حُدّ $n\,u_{2n}$
بشريحة $\sum_{k=n+1}^{2n} u_k$ واستعمل محك
كوشي)}. وبيّن أن العكس يفشل، وأن فرض
الرتابة لا يمكن إزالته.
\end{exercise}

\section{مسألة: ثابت أويلر والمتسلسلة التي تغيّر
مجموعها}

\begin{problem}[{مسألة نهاية الأسبوع --- $H_n = \ln n + \gamma +
\frac{1}{2n} + O(n^{-2})$، وإعادة ترتيب $1 - \frac12 +
\frac13 - \dots$ إلى $\frac{\ln 2}{2}$}]
\label{pb:b1:series:1}
قصتان تتقاسمان \omterm{def:b1:series:def}{المتسلسلة} التوافقية. الأولى، مسك الدفاتر
المضبوط لتباعدها: إذ يتقارب $H_n - \ln n$ إلى \omterm{pb:b1:series:1}{ثابت أويلر}
$\gamma$ (\cref{exo:b1:series:7})، وتلطّف هذه المسألة
العبارةَ إلى قانون ذي طرفين $\frac{1}{2(n+1)}
\leq H_n - \ln n - \gamma \leq \frac{1}{2n}$، مصادِقةً على $\gamma
= 0.5772\dots$ باليد. والثانية، فضيحة التقارب
الشرطي: تجمع \omterm{def:b1:series:def}{المتسلسلة} التوافقية المتناوبة إلى $\ln 2$
(\cref{ex:b1:series:ln2})، ومع ذلك تجمع \emph{الحدود نفسها، بترتيب
مختلف}، إلى $\frac{\ln 2}{2}$ --- أو إلى $\ln 2 +
\frac12\ln\frac pq$ من أجل أيّ $p, q$، أو إلى أيّ عدد حقيقي كان
(ريمان). والقصتان قصة واحدة: إذ تُحسب المجاميع المعاد ترتيبها
\emph{بقانون} $\gamma$.\index{ثابت أويلر}\index{إعادة ترتيب المتسلسلات}

\medskip
\noindent\textbf{الجزء 1 --- $\gamma$، مؤطَّرًا.} ضع $a_n = H_n
- \ln n$ و $b_n = H_n - \ln(n+1)$.
\begin{enumerate}
  \item باستعمال $\frac{t}{1+t} \leq \ln(1 + t) \leq t$، بيّن أن
        $(a_n)$ متناقصة و $(b_n)$ متزايدة، وأنهما
        متجاورتان؛ ونهايتهما المشتركة هي $\gamma$، مع $b_n \leq
        \gamma \leq a_n$ من أجل كل $n$.
  \item أول طلقة عددية: انطلاقًا من $H_{10} = 2.928968\dots$،
        أطّر $\gamma$ بين $b_{10} = 0.5311$ و $a_{10} =
        0.6264$. وكم يجب أن يكبر $n$ في هذا التأطير الفجّ
        من أجل أربعة أرقام عشرية؟
  \item بيّن التمثيل الذيلي المضبوط $a_n - \gamma =
        \sum_{k \geq n} w_k$ (نهاية المجاميع الجزئية)، حيث
        \[
        w_k = a_k - a_{k+1}
        = \ln\Bigl(1 + \frac1k\Bigr) - \frac{1}{k+1}
        = \int_k^{k+1} \frac{(k + 1 - t)}{t\,(k+1)}\,\dd t ,
        \]
        واستنتج من الصورة التكاملية الحاصر ذا الطرفين
        $\dfrac{1}{2(k+1)^2} \leq w_k \leq \dfrac{1}{2k(k+1)}$.
\end{enumerate}

\medskip
\noindent\textbf{الجزء 2 --- قانون $\frac{1}{2n}$.}
\begin{enumerate}[resume]
  \item اجمع حواصر السؤال 3 (فالطرفان يتلسكبان أو
        يُقارنان بتلسكوبات) واستنتج القانون:
        \[
        \frac{1}{2(n+1)} \;\leq\; H_n - \ln n - \gamma \;\leq\;
        \frac{1}{2n} \qquad (n \geq 1).
        \]
  \item استنتج $H_n = \ln n + \gamma + \frac{1}{2n} +
        O\bigl(\frac{1}{n^2}\bigr)$؛ وبدقة، بيّن أن
        $\gamma_n = H_n - \ln n - \frac{1}{2n}$ يحقق
        $-\frac{1}{2n(n+1)} \leq \gamma_n - \gamma \leq 0$.
  \item صادِق على أربعة أرقام عشرية مع $n = 100$: فمن $H_{100} =
        5.1873775\dots$، احسب $\gamma_{100} = 0.577207\dots$
        واستنتج $\gamma = 0.5772 \pm 5\cdot10^{-5}$ (والقيمة
        الحقيقية $0.5772156\dots$).
  \item ربحان من القانون، وكلاهما لازم لاحقًا: عندما $m \to
        \infty$،
        \[
        H_{2m} - H_m = \ln 2 - \frac{1}{4m} +
        O\Bigl(\frac{1}{m^2}\Bigr),
        \qquad
        \sum_{j=1}^{m} \frac{1}{2j-1} = \frac{\ln m}{2} + \ln 2
        + \frac\gamma2 + o(1) ,
        \]
        والثاني عبر $\sum_{j \leq m} \frac{1}{2j-1} = H_{2m}
        - \frac12 H_m$، وكذلك $\sum_{j=1}^{m}
        \frac{1}{2j} = \frac{\ln m}{2} + \frac\gamma2 + o(1)$.
\end{enumerate}

\medskip
\noindent\textbf{الجزء 3 --- \omterm{def:b1:series:def}{المتسلسلة} التوافقية المتناوبة،
حتى الرتبة الثانية.}
\begin{enumerate}[resume]
  \item بيّن (بالاستقراء أو بالتجميع) المتطابقة
        $\sum_{k=1}^{2m} \frac{(-1)^{k-1}}{k} = H_{2m} - H_m$،
        واستنتج المجموع $\ln 2$ (مرة أخرى) والسرعة
        المضبوطة:
        \[
        \sum_{k=1}^{2m} \frac{(-1)^{k-1}}{k}
        = \ln 2 - \frac{1}{4m} + O\Bigl(\frac{1}{m^2}\Bigr) .
        \]
  \item استنتج الخطأ المقارب \omterm{def:b1:series:def}{للمتسلسلة} التوافقية
        المتناوبة عند \emph{أيّ} دليل: $S - S_N \sim
        \frac{(-1)^N}{2N}$ --- أي أصغر مرتين من حاصر
        أسوأ حالة $a_{N+1} \approx \frac1N$ في
        \cref{thm:b1:series:alternating}.
  \item (تسريع بالمجّان) بيّن أن المجاميع الممتوسطة
        $\tilde S_N = \frac{S_N + S_{N+1}}{2}$ تحقق $\tilde
        S_N = \ln 2 + O\bigl(\frac{1}{N^2}\bigr)$. وللتحقق: $S_{10}
        = 0.64563$ و $S_{11} = 0.73654$ و $\tilde S_{10} =
        0.69109$، إزاء $\ln 2 = 0.69315$: فمتوسط واحد يشتري
        منزلتين عشريتين.
  \item فسّر في جملتين لماذا لا يمكن لمثل هذه الحيلة أن تساعد ظاهرةً
        ذات ذيل \emph{موجب} متباعد مثل تأطير السؤال
        2: فخطأ المتناوب يتذبذب (بإشارة
        $(-1)^N$)، ومنه يمحو المتوسط حدَّه المهيمن، بينما
        لخطأ تأطير $\gamma$ وهو $\frac{1}{2n}$ إشارةٌ
        ثابتة. (وأمّا متوسط $a_n$ و $b_n$ \emph{يساعد} فعلًا:
        اربط $\frac{a_n + b_n}{2}$ بتقدير المنتصف
        $H_n - \ln\bigl(n + \frac12\bigr)$ وبيّن أن خطأه هو
        $O\bigl(\frac{1}{n^2}\bigr)$.)
\end{enumerate}

\medskip
\noindent\textbf{الجزء 4 --- الصلابة وفشلها.}
\begin{enumerate}[resume]
  \item بيّن أن الجزء الموجب $\sum \frac{1}{2j-1}$ و
        الجزء السالب $\sum \frac{1}{2j}$ من \omterm{def:b1:series:def}{المتسلسلة} التوافقية
        المتناوبة كليهما يتباعد --- وهو توقيع
        التقارب \emph{الشرطي}.
  \item برهن على \omterm{def:b1:logic:statement}{العبارة} العامة وراء السؤال 12: إذا
        تقاربت $\sum u_n$ وتباعدت $\sum \abs{u_n}$، فإن
        \omterm{def:b1:series:def}{متسلسلة} الأجزاء الموجبة $\sum u_n^+$ \omterm{def:b1:series:def}{ومتسلسلة} الأجزاء
        السالبة $\sum u_n^-$ كلتيهما تتباعد \emph{(من
        $u_n^\pm = \frac{\abs{u_n} \pm u_n}{2}$: فلو تقاربت
        إحداهما، لتقاربت الأخرى، ومنه $\sum\abs{u_n}$)}.
        وهذا الخزان الذي لا ينضب من الكتلة الموجبة والسالبة
        هو ما ستنفقه وصفة ريمان.
  \item (الصلابة) برهن: إذا تقاربت $\sum u_n$
        \emph{بإطلاق} وكان $\sigma \colon \N \to \N$
        تقابلًا، فإن $\sum u_{\sigma(n)}$ تتقارب إلى
        المجموع نفسه \emph{(فمن أجل $N$ كبير تحتوي الحدود $M$ الأولى
        المعاد ترتيبها $u_0, \dots, u_N$؛ فقارن المجاميع الجزئية
        عبر الذيل $\sum_{n > N}\abs{u_n}$)}.
  \item (وصفة ريمان) ليكن $t \in \R$. صف إعادة الترتيب
        الجشعة \omterm{def:b1:series:def}{للمتسلسلة} التوافقية المتناوبة: خذ
        حدودًا موجبة $1, \frac13, \frac15, \dots$ حتى يتجاوز
        المجموع الجزئي $t$ أول مرة، ثم حدودًا سالبة
        $-\frac12, -\frac14, \dots$ حتى يهبط دون
        $t$ أول مرة، وكرّر. بيّن أن كل حدّ يُستعمل مرة واحدة
        بالضبط، وأنه بعد أول عبور تبقى المجاميع الجزئية
        في حدود آخر حدّ مستعمل من $t$، واستنتج أن
        \omterm{def:b1:series:def}{المتسلسلة} المعاد ترتيبها تتقارب إلى $t$: أي إن \emph{أيّ}
        مجموع مفروض قابل للبلوغ.
\end{enumerate}

\medskip
\noindent\textbf{الجزء 5 --- صيغة $(p, q)$.} ثبّت عددين صحيحين
$p, q \geq 1$. أعد ترتيب \omterm{def:b1:series:def}{المتسلسلة} التوافقية المتناوبة في
كتل: $p$ حدًّا موجبًا (وهي مقلوبات الأعداد الفردية التالية)، ثم $q$
حدًّا سالبًا (وهي مقلوبات الأعداد الزوجية التالية)، وكرّر.
\begin{enumerate}[resume]
  \item تحقق من أن هذه إعادة ترتيب حقيقية (كل حدّ
        مرة واحدة بالضبط)، وأنها من أجل $(p, q) = (1, 2)$ تُقرأ
        \[
        1 - \frac12 - \frac14 + \frac13 - \frac16 - \frac18 +
        \frac15 - \frac1{10} - \frac1{12} + \dots
        \]
  \item (التنصيف المضبوط) من أجل $(p, q) = (1, 2)$، برهن على
        متطابقة الكتلة
        \[
        \frac{1}{2k-1} - \frac{1}{4k-2} - \frac{1}{4k}
        = \frac12\Bigl(\frac{1}{2k-1} - \frac{1}{2k}\Bigr),
        \]
        واستنتج العلاقة \emph{المضبوطة} $T_{3K} = \frac12
        S_{2K}$ بين المجاميع الجزئية المعاد ترتيبها و
        الأصلية: فتنصيف المجموع مرئي عند
        كل مرحلة منتهية، لا في النهاية فقط.
  \item بيّن أن المجموع الجزئي بعد $K$ كتلة كاملة
        يساوي $\sum_{j=1}^{pK} \frac{1}{2j-1} -
        \sum_{j=1}^{qK} \frac{1}{2j}$، واحسب نهايته
        مع السؤال 7:
        \[
        \ln 2 + \frac12 \ln\frac pq .
        \]
  \item اضبط المجاميع الجزئية \emph{داخل} كتلة (فالحدود تؤول إلى $0$)
        واستنتج أن \omterm{def:b1:series:def}{المتسلسلة} المعاد ترتيبها بالمقدار $(p,
        q)$ تتقارب إلى $\ln 2 + \frac12\ln
        \frac pq$. وعلى الخصوص يعطي $(1, 2)$ المقدارَ $\frac{\ln
        2}{2}$: فتحقق إزاء الحدود التسعة الأولى، $T_9 =
        0.3083$، وهي تزحف نحو $0.3466$.
  \item تحققات سلامة ومدى: يستعيد $(1,1)$ المقدارَ $\ln 2$؛
        ويعطي $(2,1)$ المقدارَ $\frac32\ln 2$؛ وأيّ المجاميع يمكن بلوغها
        بكتل $(p, q)$، وكيف تقارن هذه القائمة القابلة للعدّ
        بقائمة ريمان الكاملة (السؤال 14)؟
\end{enumerate}

\medskip
\noindent\textbf{الجزء 6 --- خاتمة: $\gamma$ في العمل، و
توليفة.}
\begin{enumerate}[resume]
  \item حدّد مجموع \omterm{def:b1:series:def}{المتسلسلة} المتقاربة
        $\sum_{k\geq1} \bigl(\frac1k - \ln\frac{k+1}{k}\bigr)$
        (\cref{exo:b1:series:7}): بيّن أنه يساوي $\gamma$.
  \item نفّذ وصفة ريمان (السؤال 14) من أجل الهدف $t =
        1$ واذكر الحدود الاثني عشر الأولى الناتجة ($1, \frac13,
        -\frac12, \frac15, -\frac14, \frac17, \frac19, -\frac16,
        \frac1{11}, \frac1{13}, -\frac18, \frac1{15}$)،
        حاسبًا المجموع الجزئي ($\approx 0.980$) --- وراقب
        الخوارزمية تتنفس حول هدفها.
  \item بيّن أن إعادة ترتيب ما \omterm{def:b1:series:def}{للمتسلسلة} التوافقية المتناوبة
        تتباعد إلى $+\infty$ \emph{(بكتل من الحدود
        الموجبة طويلة بما يكفي لكسب $1$ في كل مرة، باستعمال السؤال
        12، يفصل بينها حدود سالبة مفردة)}.
  \item لطّف \cref{ex:b1:series:harmonicstack} بقانون
        $\gamma$: بيّن أن أول دليل يحقق $H_N \geq
        20$ يحقق $N = \eu^{\,20 - \gamma}\,(1 + o(1))
        \approx 2.7\cdot10^{8}$ --- \omterm{pb:b1:series:1}{فثابت أويلر} هو بالضبط
        التصحيح الذي كان ينقص التأطيرَ الفجّ.
  \item توليفة، جملة واحدة لكلٍّ: (أ) قانون $\gamma$ و
        ما يسهم به كلٌّ من قطعه الثلاث ($\ln n$ و $\gamma$ و
        $\frac{1}{2n}$)؛ (ب) ولماذا يجعل التقارب الشرطي
        المجموعَ متعلقًا بالترتيب بينما يمنع \omterm{thm:b1:series:absolute}{التقارب
        بإطلاق} ذلك؛ (ج) وكيف كانت صيغة $(p,q)$
        \emph{حسابًا} بقانون $\gamma$ لا ادعاءَ
        وجود مجردًا؛ (د) وأين تستمر هذه
        الخيوط --- المتسلسلات الصحيحة وجداءات المتسلسلات
        في مجلّد السنة الثانية، ومسألة نهاية الأسبوع في مجلّد السنة الثالثة
        عن صيغة ستيرلينغ، حيث يجري مسك الدفاتر نفسه
        بين المجاميع والتكاملات بكامل قوته.
\end{enumerate}
\end{problem}
